\documentclass[../main.tex]{subfiles}

\begin{document}

% -----------------------------
\section{Parsing mlir file}
% -----------------------------
Now let's parse and print created mlir file from Python:
Create a file \path{step1_parse_and_print.py}

\begin{lstlisting}[language=Python, caption={step1\_parse\_and\_print.py}]
from xdsl.context import Context
from xdsl.parser import Parser
from xdsl.dialects.builtin import Builtin
from xdsl.dialects.func import Func
from xdsl.dialects.arith import Arith

def main() -> None:
    # Create context and load dialects you need
    ctx = Context()
    ctx.load_dialect(Builtin)
    ctx.load_dialect(Func)
    ctx.load_dialect(Arith)

    # Read MLIR/xDSL file
    with open("test.mlir", "r") as f:
        text = f.read()

    # Parse into a module
    parser = Parser(ctx, text)
    module = parser.parse_module()

    # Print the module back out
    print("=== Parsed module ===")
    print(module)


if __name__ == "__main__":
    main()
\end{lstlisting}

\begin{lstlisting} [language=bash, caption={Parse file}]
python3 step1_parse_and_print.py
\end{lstlisting}

\textbf{Exercises}
\begin{itemize}
    \item Create a mlir file for \path{mul_two} function (Hint: use arith.muli)
    \item Create a mlir file for \path{sub_two} function (Hint: use arith.subi)
    \item Create a mlir file for linear function \path{ax + b} (3 arguments)
\end{itemize}
\end{document}
